\section{The Hungerbühler--Wasem generalised residue theorem}\label{sec:hw33}

In this chapter we present the dependency chain culminating in the
fully discharged paper-faithful form of the Hungerbühler--Wasem
generalised residue theorem (Hungerbühler--Wasem 2018, Theorem 3.3).
The headline statement asserts that the multi-point Cauchy principal
value of $f$ along a closed null-homologous piecewise-$C^{1}$
immersion equals a weighted sum of residues, where the weights are
the generalised winding numbers. The proof proceeds by reducing to
the per-pole problem, splitting each pole into its simple and
higher-order Laurent contributions, and handling the higher-order
part by the sector-cancellation identity under condition (B).

\begin{theorem}[HW3.3 — paper-faithful clean form]
  \label{thm:residue-paper-faithful-clean}
  \lean{HungerbuhlerWasem.residueTheorem_crossing_paper_faithful_clean}
  \uses{def:closed-pwc1-immersion, def:null-homologous, def:has-cauchy-pv-on,
    def:generalized-winding-number, def:condition-a-prime, def:condition-b}
  \leanok
  Let $U \subseteq \mathbb{C}$ be a non-empty open set, $S \subseteq U$
  a finite set, and $f : \mathbb{C} \to \mathbb{C}$ a function. Assume:
  \begin{itemize}
    \item $f$ is holomorphic on $U \setminus S$ and meromorphic at
      every point of $S$;
    \item $\gamma$ is a closed piecewise-$C^{1}$ immersion based at
      $x \in U$, with $x \notin S$;
    \item $\gamma$ is null-homologous in $U$;
    \item $(\gamma, f)$ satisfies condition (B) relative to $S$;
    \item $\gamma$ satisfies condition (A') relative to $S$ with
      $\operatorname{poleOrder}(s)$ equal to the order of the polar
      part of $f$ at $s$ provided by the canonical polar-part
      decomposition.
  \end{itemize}
  Then the multi-point Cauchy principal value of $\oint_\gamma f$
  exists and equals
  \[
    \operatorname{PV}_S\!\!\int_\gamma f(z)\,\mathrm{d}z
    \;=\; \sum_{s \in S}
      2 \pi i \, \operatorname{gWN}(\gamma, s) \, \operatorname{Res}(f, s).
  \]
\end{theorem}

\begin{proof}
  \uses{thm:residue-compositional, thm:cpv-polar-uncrossed,
    thm:cpv-polar-multicrossed-condB}
  \leanok
  By extracting the canonical polar-part decomposition of $f$ over
  $S$ in $U$ from the condition-(B) data, we reduce to the
  compositional residue theorem (\cref{thm:residue-compositional}).
  That reduction in turn requires, for each $s \in S$, an explicit
  per-pole Cauchy principal value identity for the polar part. We
  split the proof of those identities into two cases:
  \begin{itemize}
    \item if $\gamma$ avoids the pole $s$ on $[0,1]$, the polar part
      is regular along $\gamma$, and the principal value coincides
      with the ordinary contour integral; this is
      \cref{thm:cpv-polar-uncrossed};
    \item otherwise the set of crossings of $\gamma$ at $s$ is a
      finite subset of $(0,1)$ (its finiteness uses
      $x \notin S$); we then combine the corner-friendly
      multi-crossing CPV identity \cref{thm:cpv-polar-multicrossed-condB}
      with the canonical right/left tangent limits at each crossing.
  \end{itemize}
  Together these per-pole identities, applied to every $s \in S$,
  combine via the compositional theorem to yield the stated
  equality.
\end{proof}

\begin{theorem}[Compositional residue theorem]
  \label{thm:residue-compositional}
  \lean{HungerbuhlerWasem.residueTheorem_crossing_compositional}
  \uses{def:closed-pwc1-immersion, def:null-homologous, def:has-cauchy-pv-on,
    def:generalized-winding-number}
  \leanok
  Let $U \subseteq \mathbb{C}$ be a non-empty open set,
  $S \subseteq U$ a finite set, $f : \mathbb{C} \to \mathbb{C}$
  holomorphic on $U \setminus S$, and $\gamma$ a closed piecewise-$C^{1}$
  immersion based at $x \in U$ that is null-homologous in $U$. Assume
  that $f$ admits a polar-part decomposition
  $f(z) = h(z) + \sum_{s \in S} P_s(z)$, where each $P_s$ is the
  Laurent polar part at $s$ and $h$ is the analytic remainder, and
  that for each $s \in S$ the polar part has principal value
  \[
    \operatorname{PV}_S\!\!\int_\gamma P_s(z)\,\mathrm{d}z
      \;=\; 2 \pi i \, \operatorname{gWN}(\gamma, s) \,
        \operatorname{Res}(f, s).
  \]
  Then
  \[
    \operatorname{PV}_S\!\!\int_\gamma f(z)\,\mathrm{d}z
    \;=\; \sum_{s \in S}
      2 \pi i \, \operatorname{gWN}(\gamma, s) \,
        \operatorname{Res}(f, s).
  \]
\end{theorem}

\begin{proof}
  \leanok
  The analytic remainder $h$ has principal value zero by Dixon's
  theorem applied to the null-homologous $\gamma$ together with the
  observation that the multi-point CPV-cutoff integral of an analytic
  function on $U$ along a null-homologous curve tends to the contour
  integral of $h$, which vanishes. Adding this to the assumed
  per-pole identities, and exchanging finite sum with the
  $\varepsilon \to 0^{+}$ limit, gives the result. The
  exchange-of-limit step is justified by the integrability of each
  cutoff integrand on $[0,1]$, which holds because the polar parts
  are continuous off their poles and $h$ is continuous everywhere.
\end{proof}

\begin{theorem}[Per-pole CPV at an uncrossed pole]
  \label{thm:cpv-polar-uncrossed}
  \lean{HungerbuhlerWasem.cpv_polarPart_at_uncrossed_pole}
  \uses{def:closed-pwc1-immersion, def:null-homologous, def:has-cauchy-pv-on,
    def:generalized-winding-number}
  \leanok
  Let $U \subseteq \mathbb{C}$ be open, $S \subseteq U$ finite, and
  $f : \mathbb{C} \to \mathbb{C}$ with polar-part decomposition over
  $S$ in $U$. Let $\gamma$ be a closed piecewise-$C^{1}$
  immersion null-homologous in $U$, and let $s \in S$ be a pole that
  $\gamma$ \emph{avoids}: $\gamma(t) \ne s$ for every $t \in [0,1]$.
  Then the polar part at $s$ has multi-point Cauchy principal value
  \[
    \operatorname{PV}_S\!\!\int_\gamma P_s(z)\,\mathrm{d}z
    \;=\; 2 \pi i \, \operatorname{gWN}(\gamma, s) \,
      \operatorname{Res}(f, s).
  \]
\end{theorem}

\begin{proof}
  \leanok
  Since $\gamma$ avoids $s$, the distance from $\gamma$ to $s$ is
  bounded below by some $\delta > 0$ by compactness of the image.
  For $\varepsilon < \delta$ the multi-point cutoff integrand for the
  polar part coincides with the ordinary integrand on all of $[0,1]$,
  so the cutoff integrals are eventually constant in $\varepsilon$
  and equal to the contour integral
  $\int_\gamma P_s(z)\,\mathrm{d}z$. The dominated-convergence
  argument therefore reduces to evaluating the contour integral of
  $P_s$ on $\gamma$, and the avoidance computation identifies this
  with $2\pi i \, \operatorname{gWN}(\gamma, s) \, \operatorname{Res}(f, s)$.
\end{proof}

\begin{theorem}[Per-pole CPV at a multi-crossed pole under condition (B)]
  \label{thm:cpv-polar-multicrossed-condB}
  \lean{HungerbuhlerWasem.cpv_polarPart_at_multiCrossed_pole_under_condB_corner}
  \uses{def:closed-pwc1-immersion, def:has-cauchy-pv-on,
    def:generalized-winding-number, def:condition-a-prime, def:condition-b}
  \leanok
  Let $\gamma$ be a closed piecewise-$C^{1}$ immersion at $x$, let
  $s \in S$ be a pole at which the polar part of $f$ has order $N$,
  and let
  \[
    \operatorname{crossings} = \{ t \in (0,1) : \gamma(t) = s \}
  \]
  be a finite (corner-allowing) crossing set. Assume the following
  data on $\operatorname{crossings}$:
  \begin{itemize}
    \item at every $t \in \operatorname{crossings}$ the curve
      $\gamma$ is flat of order $N$ (the condition (A') witness);
    \item right and left tangent limits $L_+(t), L_-(t) \in
      \mathbb{C}^{\times}$ exist;
    \item the angle-compatibility identity from condition (B)
      \[
        \bigl(L_+(t) / |L_+(t)|\bigr)^{k}
        = \bigl(-L_-(t) / |L_-(t)|\bigr)^{k}
      \]
      holds for every $k \ge 1$ at which the Laurent coefficient
      $a_k$ at $s$ is non-zero;
    \item the simple-pole Cauchy principal value of $(z - s)^{-1}$
      along $\gamma$ exists.
  \end{itemize}
  Then the polar part at $s$ has multi-point CPV
  \[
    \operatorname{PV}_{\{s\}}\!\!\int_\gamma P_s(z)\,\mathrm{d}z
    \;=\; 2 \pi i \, \operatorname{gWN}(\gamma, s) \,
      \operatorname{Res}(f, s).
  \]
\end{theorem}

\begin{proof}
  \leanok
  We split the polar part as
  $P_s(z) = a_0/(z-s) + \sum_{k=1}^{N-1} a_k/(z-s)^{k+1}$. The
  $k = 0$ term is the simple-pole contribution: by the assumed
  existence of the principal value for $(z-s)^{-1}$, that principal
  value equals $2\pi i \, \operatorname{gWN}(\gamma, s)$, so the
  $k=0$ contribution gives $a_0$ times this expression, which equals
  $2\pi i \, \operatorname{gWN}(\gamma, s) \, \operatorname{Res}(f, s)$.

  For each $k \ge 1$, two cases. If $a_k = 0$ the term is
  identically zero. If $a_k \ne 0$, the flatness-of-order-$N$
  hypothesis and the angle-compatibility condition (B) hypothesis
  combine to give the sector-cancellation identity: the cutoff
  integral of $a_k/(z-s)^{k+1}$ along $\gamma$ vanishes in the
  $\varepsilon \to 0^+$ limit. Summing over $k$ from $0$ to $N-1$
  yields the claim.
\end{proof}

\begin{theorem}[HW3.3 -- fully general, multi-crossing, meromorphic form]
  \label{thm:hw-3-3-clean-full-mero}
  \lean{LeanModularForms.hw_3_3_clean_full_mero}
  \uses{thm:residue-paper-faithful-clean, def:closed-pwc1-immersion,
    def:null-homologous, def:has-cauchy-pv-on, def:generalized-winding-number,
    def:condition-a-prime, def:condition-b}
  \leanok
  Let $U \subseteq \mathbb{C}$ be a non-empty open set, $S \subseteq U$
  a finite set, and $f : \mathbb{C} \to \mathbb{C}$ a function such
  that
  \begin{itemize}
    \item $f$ is holomorphic on $U \setminus S$, and
    \item $f$ is meromorphic at each $s \in S$.
  \end{itemize}
  Let $\gamma$ be a closed piecewise-$C^{1}$ immersion based at
  $x \in U$ with $x \notin S$, null-homologous in $U$. Suppose
  that $(\gamma, f)$ satisfies condition (B) and that $\gamma$
  satisfies condition (A') with pole-order data equal to the
  canonical polar-part orders extracted from the
  condition-(B) witness. Then
  \[
    \operatorname{PV}_S\!\!\int_\gamma f(z)\,\mathrm{d}z
    \;=\; \sum_{s \in S} 2 \pi i \,
      \operatorname{gWN}(\gamma, s) \, \operatorname{Res}(f, s).
  \]
  This is the maximally general paper-faithful form of
  Hungerbühler--Wasem 2018, Theorem 3.3, with all six oracle
  hypotheses of the intermediate ``\emph{hw\_3\_3\_paper}'' formulation
  discharged.
\end{theorem}

\begin{proof}
  \uses{thm:residue-paper-faithful-clean}
  \leanok
  This is a direct invocation of the paper-faithful clean form
  \cref{thm:residue-paper-faithful-clean}: the hypothesis list of
  the present theorem matches that of the clean form term-for-term,
  and no additional reduction is needed.

  The single structural residual $x \notin S$ is the only
  Lean-formalisation artefact (the curve in our setting carries a
  typed basepoint, whereas the paper works with cycles). It is
  routinely satisfiable: if $x \in S$, replace $\gamma$ by the
  cyclic shift $\gamma(\,\cdot\, + \tau)$ for a parameter
  $\tau \in (0,1)$ at which $\gamma(\tau) \notin S$, which exists
  since the preimage of the finite set $S$ under $\gamma$ is
  finite.
\end{proof}
